% !TEX root = task-2-6-7.tex
% Editable task bank. Append new BankTask blocks at the end of the document.
\documentclass[12pt,a4paper]{article}
\usepackage{physics-textbook}
\usepackage{graphicx,needspace}
\usetikzlibrary{decorations.pathmorphing}
\newif\ifShowSolutions
\ifdefined\TutorEdition
  \ShowSolutionstrue
\else
  \ShowSolutionsfalse % Change to \ShowSolutionstrue to display solutions.
\fi
\newcounter{banktask}
\newsavebox{\BankTaskBox}
\newcommand{\BankTaskContent}[3]{%
  \noindent\textbf{\thebanktask.}\quad #2\par
  \ifShowSolutions
    \Needspace{50mm}% Keep the solution heading with its opening text or diagram.
    #3\par
    \ifdefined\TutorEdition
      \ifcsname TutorSource#1\endcsname
        \nopagebreak\smallskip{\small\csname TutorSource#1\endcsname\par}%
      \fi
    \fi
  \fi
}
\newcommand{\BankTask}[3]{%
  \refstepcounter{banktask}\par\medskip
  \sbox{\BankTaskBox}{\begin{minipage}{\linewidth}%
    \BankTaskContent{#1}{#2}{#3}\end{minipage}}%
  \ifdim\dimexpr\ht\BankTaskBox+\dp\BankTaskBox\relax>\dimexpr\textheight-20mm\relax
    \Needspace{35mm}\BankTaskContent{#1}{#2}{#3}%
  \else
    \Needspace{\dimexpr\ht\BankTaskBox+\dp\BankTaskBox\relax}%
    \noindent\usebox{\BankTaskBox}\par
  \fi
  \medskip
}
\begin{document}
\BookSubsection{2.6.7}{Графики изохорного процесса}


% Source 79B504: https://ege.fipi.ru/bank/questions.php?search=1&proj=BA1F39653304A5B041B656915DC36B38&qid=79B504
\BankTask{79B504}{
\begin{center}
\begin{tikzpicture}[x=0.11cm,y=5.5cm,>=stealth]
\draw[thin,gray] (0,0) grid[xstep=10,ystep=0.2] (40,0.6);
\draw[->] (0,0)--(51,0) node[right] {$T,\,\text{К}$};
\draw[->] (0,0)--(0,0.82) node[left] {$p,\,10^{4}\,\text{Па}$};
\foreach \x/\lab in {0/290,10/300,20/310,30/320,40/330} {
\draw (\x,0)--(\x,-0.015);
\node[below] at (\x,-0.015) {$\lab$};
}
\foreach \y/\lab in {0/{1{,}8},0.2/{2{,}0},0.4/{2{,}2},0.6/{2{,}4}} {
\draw (-0.7,\y)--(0,\y);
\node[left] at (-0.7,\y) {$\lab$};
}
\draw[thick] (0,0.1333)--(44,0.4267);
\end{tikzpicture}
\end{center}
На рисунке показан график зависимости давления газа в запаянном сосуде от его температуры. Объём сосуда равен $0{,}25\,\text{м}^{3}$. Какое приблизительно количество газообразного вещества содержится в этом сосуде (в $\text{моль}$, с округлением до целого)?
}{
\smallskip\noindent\textbf{Решение от GPT}\par
По графику при $T=300\,\text{К}$ давление $p\approx2{,}0\cdot10^{4}\,\text{Па}$. Из уравнения идеального газа
\[
\nu=\frac{pV}{RT}=\frac{2{,}0\cdot10^{4}\cdot0{,}25}{8{,}31\cdot300}\approx2{,}01\,\text{моль}.
\]
После округления ответ: $2\,\text{моль}$.
}

\end{document}
